\section{$\cN{=}2$ 理论的谱系}\label{sec:zoo}
在本讲义中，我们见到了许多四维 $\cN{=}2$ 理论。本附录旨在将它们集中汇总，因为其中许多理论是在行文过程中逐渐引入的，分散在不同位置。\footnote{作者感谢 Simone Giacomelli 帮助准备本节。}

遗憾的是，由于对全部 $\cN{=}2$ 理论的理解仍很不充分，我们无法对涉及的所有理论给出一个连贯的分类。不过，已经能够看出一些规律：存在许多标准的无穷系列，以及一些“例外”理论。此外，无穷系列最前面的某些成员会偶然地彼此相同。这与单李代数的分类如出一辙：存在四个无穷系列 $A_k=\SU(k{+}1)$、$B_k=\SO(2k{+}1)$、$C_k=\Sp(k)$ 和 $D_k=\SO(2k)$，以及例外的 $E_{6,7,8}$、$F_4$ 和 $G_2$。其中有 $A_1=B_1=C_1$、$B_2=C_2$、$A_3=D_3$ 等偶然的重合，而且某些例外系列可以自然地向下延伸，使得 $E_5=D_5$、$E_4=A_4$。

\subsection{规范理论}\label{sec:zoogauge}\index{QuiverGaugeTheory@箭图规范理论|hyperpage}\index{SemiLagrangianTheory@半拉格朗日理论|hyperpage}
首先讨论拉格朗日规范理论。正如第 \ref{sec:lagrangian} 节所详细讨论的，一个拉格朗日 $\cN{=}2$ 超对称理论基本上由以下内容指定：
\begin{itemize}
\item 一个规范群 $G$，它可以是许多单因子的乘积；
\item 以及一个属于 $G$ 的表示 $R$ 的（半）超多重态。
\end{itemize}
此外，规范群 $G$ 的每个单因子都有一个复化规范耦合，并且可以为超多重态添加质量项。

为了使理论在紫外完备，即在每个能标上都有定义，式 \eqref{n2running} 所定义的单圈β 函数必须非负。对给定的 $G$，这对可能的超多重态表示 $R$ 施加了约束。对于单群 $G$，所有这类允许表示的列表早已为人所知 \cite{Koh:1983ir,Dong:1984vt,Derendinger:1984bu}。半单群 $G$ 的分类在文献 \cite{Bhardwaj:2013qia} 中完成。

一个规范群 $G$ 为半单群的典型规范理论可按如下方式获得。取一个由顶点和边组成的图。对每个顶点 $v$，引入一个 $\SU(N_v)$ 规范多重态；对连接顶点 $v$ 和 $v'$ 的每条边 $e$，引入一个 $\SU(N_v)\times\SU(N_{v'})$ 的双基本表示超多重态。给定一个图，可以找到颜色数 $N_v$ 的一组取值，使得
\begin{itemize}
\item 理论渐近自由，当且仅当该图是有限型 Dynkin 图 $\Gamma$；
\item 理论超共形，当且仅当该图是仿射 Dynkin 图 $\hat\Gamma$。
\end{itemize}
它们通常分别称为 $\Gamma$ 型或 $\hat\Gamma$ 型箭图规范理论，最初在文献 \cite{Katz:1997eq} 中得到系统分析。

更一般地，可以从一个具有味对称性 $F$、可能是非拉格朗日的理论 $T$ 出发。然后选取子群 $G\subset F$，并尝试将规范群为 $G$ 的矢量多重态耦合到理论 $T$。只要按式 \eqref{GOO} 定义的 $G$ 的流的两点函数系数不太大，组合后的系统就应在每个能标上都有意义。这样的组合系统可称为半拉格朗日系统，因为其中既有尚无拉格朗日描述的部分 $T$，也有具有拉格朗日描述的矢量多重态部分。下面不明确列出这类半拉格朗日系统。

\subsection{六维构造}\label{sec:app6d}
在本讲义中，拉格朗日规范理论的 Seiberg-Witten 解，几乎总是用某个紧致化在紫外曲线上的六维 $\cN{=}(2,0)$ 超对称理论来表述，如第 \ref{sec:6d} 节所解释的那样。我们主要关注由 $N$ 个 M5 膜得到的六维理论：即 $\SU(N)=A_{N-1}$ 型六维 $\cN{=}2$ 理论。此外，已知还存在 $\SO(2N)=D_N$ 型以及 $E_6$、$E_7$ 和 $E_8$ 型六维 $\cN{=}(2,0)$ 理论。

\subsubsection{六维 $\cN{=}(2,0)$ 理论的类型}
对于 $A_1$ 型理论，我们在第 \ref{sec:gaiotto} 节中见到了六维世界体积上的场 $\phi(z)$，Seiberg-Witten 曲线由 $\lambda^2=\phi(z)$ 给出。在那里还看到，可以在紫外曲线上引入各种穿孔，使场 $\phi(z)$ 具有各种阶数的极点。对于 $A_{N-1}$ 型理论，我们在第 \ref{sec:duality} 节中见到了场 $\phi_2(z),\ldots,\phi_N(z)$。同样，在那里讨论了紫外曲线上的各种穿孔，它们为场 $\phi_{2,\ldots,N}(z)$ 引入各种阶数的极点。注意，$A_1$ 型的场 $\phi(z)$，在更一般的记号中对应于 $\phi_2(z)$。

我们在第 \ref{sec:SO} 小节中隐含地讨论了 $D_N$ 型六维理论。在那里，使用了世界体积场 $\phi_2(z)$、$\phi_4(z)$、$\ldots$、$\phi_{2N-2}(z)$ 和 $\phi_{2N}(z)$。不过，已知最后一个场应视为另一个场的平方：$\phi_{2N}(z)=\tilde\phi_N(z)^2$。对于 $E_n$ 型理论，存在世界体积场 $\phi_i(z)$，其中 $i$ 遍历 $E_n$ 型 Weyl 群不变量的次数。特别地，对于 $E_6$、$E_7$、$E_8$，最大的 $i$ 分别为 $12$、$18$、$30$。注意，对任意 $G$ 型六维理论，具有最大 $i$ 的世界体积场 $\phi_i(z)$ 满足 $i=h^\vee(G)$，其中 $h^\vee$ 为对偶 Coxeter 数。

\subsubsection{六维理论的穿孔}\index{TamePuncture@驯服穿孔|hyperpage}\index{WildPuncture@野性穿孔|hyperpage}\index{NilpotentElement@幂零元素|hyperpage}
驯服穿孔是指 $\phi_i(z)$ 的极点阶数小于或等于 $i$ 的穿孔。野性穿孔是指至少一个 $\phi_i(z)$ 的极点阶数大于 $i$ 的穿孔。驯服／野性穿孔也称为正则／非正则穿孔。任意类型理论中可能的驯服穿孔类型都已得到分类，见文献 \cite{Chacaltana:2012zy}。正如第 \ref{sec:tame} 小节所讨论的，对于 $A_{N-1}$ 型，驯服穿孔由 $N$ 的一个分拆标记。将一个分拆 $(t_1,\ldots,t_k)$ 与由该分拆指定 Jordan 块分解 \eqref{nilpotent-vev} 的幂零矩阵联系起来。对于任意 $G$ 型六维理论，驯服穿孔由 $G$ 的复李代数中的一个幂零元素标记，模去 $G$ 的复化伴随作用给出的共轭。对任何 $G$ 都始终存在的一个特殊幂零元素就是零元素本身。对应的穿孔称为完全穿孔，并具有相关的味对称性 $G$。对于 $G=A_{N-1}$，第 \ref{sec:tame} 小节已对此作了详细解释。

对可能的野性穿孔类型，研究还不够充分，尽管已有关于 $A_k$ 型理论的工作 \cite{Xie:2012hs}。已知对任意类型都存在一种特殊的野性穿孔，由 $i=h^\vee(G)$ 时 $\phi_i(z)$ 的极点阶数 $N$ 指定。称其为 $N$ 阶标准野性穿孔。按定义，$N>h^\vee(G)$。对于 $G=A_1$，这正是第 \ref{sec:generalAD} 小节中较为详细讨论的穿孔类型。在该小节的图中，我们用罗马数字标注极点阶数来表示这些穿孔。

\subsubsection{基本构件}\index{ClassSTheory@S 类理论|hyperpage}\index{Tinkertoy@拼装模型|hyperpage}\index{XNGTheory@$X_N(G)$ 理论|hyperpage}\index{YNGTheory@$Y_N(G)$ 理论|hyperpage}
有了这些准备，就可以按照以下步骤，构造出相当大一类四维 $\cN{=}2$ 理论：
\begin{itemize}
\item 选择六维理论的类型 $G=A_N$、$D_N$、$E_{6,7,8}$。
\item 选取放置该六维理论的紫外曲线 $C$。
\item 在 $C$ 上选取一组由各种数据标记的驯服和野性穿孔。
\item 取低能极限，使高维成分退耦。
\end{itemize}
以这种方式得到的四维理论，在文献中称为\emph{S 类理论}。

当 $C$ 的亏格较高，或穿孔较多时，可以将紫外曲线 $C$ 切开。一般而言，这一过程会在切开曲线 $C$ 的位置产生两个完全穿孔。逆过程是将一个 $G$ 矢量多重态耦合到这两个完全穿孔处两个 $G$ 味对称性的对角 $G$ 子群。对于 $G=A_{N-1}$，第 \ref{sec:duality} 节已详细描述了这一过程。

重复这一过程，就会发现任意 S 类理论都可以由规范多重态耦合以下基本成分而获得：
\begin{enumerate}
\item 带有三个驯服穿孔的球面；或
\item 带有一个驯服穿孔和一个野性穿孔的球面；或
\item 带有一个野性穿孔的球面。
\end{enumerate}
第一类中的理论称为拼装模型（tinkertoys），已在文献 \cite{Chacaltana:2010ks,Chacaltana:2011ze,Chacaltana:2012ch,Chacaltana:2013oka} 中得到广泛研究。对于 $A_{N-1}$ 型，我们在第 \ref{sec:duality} 节中讨论了其中一些例子。当三个穿孔都是完全穿孔时，理论通常记为 $T_G$；当 $G=A_{N-1}$ 时，遵循文献 \cite{Gaiotto:2009we}，也称为 $T_N$ 理论。正如第 \ref{sec:higgsing} 小节所讨论的，所有其他拼装模型都可以通过给 $T_G$ 理论赋予适当的希格斯真空期望值而获得。

对第二类和第三类理论的研究还不够充分。对于 $A_1$ 型，我们在第 \ref{sec:generalAD} 小节中以 $X_N$ 和 $Y_N$ 的名称讨论过它们。对任意类型，都存在 $X_N$ 和 $Y_N$ 的自然推广：
\begin{itemize}
\item 理论 $X_N(G)$，由 $G$ 型六维理论的一个球面得到，该球面带有一个完全穿孔和一个 $N$ 阶标准野性穿孔；
\item 理论 $Y_N(G)$，由 $G$ 型六维理论的一个球面得到，该球面仅带有一个 $N$ 阶标准野性穿孔。
\end{itemize}
这些记号在文献中完全不是标准记号。按照文献 \cite{Cecotti:2013lda} 的记号，有
\begin{equation}
X_N(G)=D_{N-h^\vee(G)}(G).\label{lda}
\end{equation}

在第 \ref{sec:pureSUN} 小节和第 \ref{sec:pureSON} 小节中，我们见到了 $G=\SU(N)$ 和 $\SO(2N)$ 时纯 $G$ 规范理论的 Seiberg-Witten 解。一般而言，当 $G=A_{N-1}$、$D_N$ 或 $E_{6,7,8}$ 时，纯 $G$ 规范理论的解已在文献 \cite{Martinec:1995by} 中求出，其结果可以用我们的语言如下表述。理论 $X_{h^\vee(G)+1}(G)$ 是具有 $G$ 对称性的空理论。取它的两个副本，并将一个 $G$ 矢量多重态耦合到它们的完全穿孔。这样就得到纯 $G$ 规范理论，其六维实现是带有两个 $h^\vee(G)+1$ 阶标准野性穿孔的球面。使两个野性奇点碰撞，就可以容易地得到这些纯理论的 Argyres-Douglas 点，正如第 \ref{sec:generalAD} 小节中对 $A_1$ 型所做的那样。两个标准野性穿孔合并为一个 $2h^\vee(G)+2$ 阶标准野性穿孔。因此得到
\begin{equation}
AD_\text{pure}(G)=Y_{2h^\vee(G)+2}(G).\label{ady}
\end{equation}
在第 \ref{sec:ADAD} 小节中，我们学到
\begin{equation}
AD_\text{pure}(\SU(N))=Y_{N+4}(A_1),\qquad
AD_\text{pure}(\SO(2N))=X_{N+2}(A_1).\label{ADDD}
\end{equation}
因此可以看出，不同 $N$ 和 $G$ 对应的 $X_N(G)$、$Y_N(G)$ 之间存在重合。

已知 $A$、$\hat A$、$D$ 和 $\hat D$ 型箭图理论可以通过这一六维构造获得。目前还不知道 $E$ 和 $\hat E$ 型箭图理论能否以这种方式获得。

\subsection{其他弦论构造}
此外还有两类利用弦论来构造 $\cN{=}2$ 理论的方法。一类是第 \ref{sec:F} 小节简要解释的 F 理论构造；另一类是在奇异 Calabi-Yau 流形上使用 IIB 型弦论。

\subsubsection{F 理论构造}\index{FTheory@F 理论|hyperpage}\index{HigherRankTheory@高秩理论|hyperpage}
F 理论是 IIB 型弦论的一个版本，其中允许轴子--伸缩子 $\tau$ 具有由 7 膜产生的非平凡 $\SL(2,\bZ)$ 单值变换。$N$ 个 D7 膜携带 $A_{N-1}$ 规范对称性，叠放在一个 O7 平面上的 $N$ 个 D7 膜携带 $D_N$ 对称性。此外，还有携带味对称性 $A_{0,1,2}$ 和 $E_{6,7,8}$ 的 7 膜。具有味对称性 $A_{0,1,2}$ 的这种额外类型的 7 膜，通常称为 $H_{0,1,2}$ 型 7 膜，以区别于通常的 D7 膜。

现在，取时空为
\begin{equation}
\bR^{1,3}\times\bC\times\bC^2,\label{Fgeom}
\end{equation}
选取一个 $G=H_{0,1,2}$ 或 $E_{6,7,8}$ 型 7 膜，将其放在
\begin{equation}
\bR^{1,3}\times\{0\}\times\bC^2.\label{Fgeomm}
\end{equation}
最后，在下述位置引入 $k$ 个 D3 膜：
\begin{equation}
\bR^{1,3}\times\{0\}\times\{0\}.
\end{equation}
将所得 $\cN{=}2$ 理论记为 $F(G,k)$。正如第 \ref{sec:rank-1SCFT} 小节中已经解释的，
\begin{equation}
AD_{N_f=n}(\SU(2))=F(H_{n-1},1)
\end{equation}
对 $n=1,2,3$ 成立，而
\begin{equation}
MN(E_n)=F(E_n,1)
\end{equation}
对 $n=6,7,8$ 成立。$k>1$ 的理论 $F(G,k)$ 是这些理论的高秩版本。理论 $F(G,k)$ 中有 $k$ 个库仑支算符，表示 $k$ 个 D3 膜沿式 \eqref{Fgeom} 中的 $\bC$ 方向运动。$F(G,k)$ 的希格斯支描述 D3 膜在 7 膜上沿式 \eqref{Fgeom} 中的 $\bC^2$ 变成瞬子的过程，因此由群 $G$ 的 $k$ 瞬子模空间给出。

这一构造还可以进一步推广：将式 \eqref{Fgeom}、\eqref{Fgeom} 中的 $\bC^2$ 方向替换为 ALE 轨形 $\bC^/\Gamma$，其可能类型已在表 \ref{tab:ADE} 中列出。这些理论尚未得到研究。

在第 \ref{sec:duality} 节、特别是第 \ref{sec:application} 小节中，我们已经讨论了 $MN(E_6)=T_3$，以及 $MN(E_{7,8})$ 也可以表示为三穿孔球面。文献 \cite{Benini:2009gi} 给出了对一般 $k$ 将 $F(E_{6,7,8},k)$ 实现为拼装模型的方法。

\subsubsection{奇异 Calabi-Yau 空间上的 IIB 型理论}\index{CalabiYauGeometry@Calabi-Yau 几何|hyperpage}\index{Conifold@锥形奇点空间|hyperpage}\index{ALEOrbifold@ALE 轨形|hyperpage}
本讲义要解释的最后一种构造如下。复维数为三的紧致 Calabi-Yau 流形上的 IIB 型弦论，给出四维 $\cN{=}2$ 超引力系统。这一构造因其与著名的镜像对称性有关而得到非常详细的研究。这里尝试使 Calabi-Yau 流形非紧致，从而使引力退耦。为了最终得到非平凡的内容，希望使 Calabi-Yau 空间具有奇点。描述复维数为三的奇异 Calabi-Yau 空间的一种简单方法，是写下含有四个复变量的方程，例如
\begin{equation}
x^2+y^2+z^2+w^2=0,\label{conifold}
\end{equation}
称为锥形奇点空间（conifold）。已知这一锥形奇点几何上的 IIB 型弦论，只产生一个自由超多重态。

这类奇异 Calabi-Yau 几何有许多类型。例如，可以用它们描述与任意 Dynkin 图相关的箭图规范理论 \cite{Katz:1997eq}。不过，这里讨论的是给出非拉格朗日理论的 Calabi-Yau 几何。

为此，先从复维数为\emph{二}而非三的奇异 Calabi-Yau 几何开始讨论。已知这些几何全部由表 \ref{tab:ADE} 中列出的 ALE 轨形 $\bC^2/\Gamma$ 穷尽。已知所有这些奇点 $\bC^2/\Gamma$ 都可以写成含三个变量的方程，见表 \ref{tab:ADEeq}。它们始终具有形式 $x^2+P_\Gamma(y,z)=0$。

\begin{table}
\[
\begin{array}{c|l@{}l}
G_\Gamma&x^2+P(y,z)&=0\\\hline
\SU(n)&x^2+y^2+z^n&=0\\
\SO(2n)&x^2+y^2z+z^{n-1}&=0\\
E_6&x^2+y^3+z^4&=0\\
E_7&x^2+y^3+yx^3&=0\\
E_8&x^2+y^3+z^5&=0
\end{array}
\]
\caption{单链李群 $G_\Gamma$ 及相应奇点 $\bC^2/\Gamma$ 的三变量方程表示。\label{tab:ADEeq}}
\end{table}

现在，用一对 $(G_\Gamma,G_{\Gamma'})$ 标记复维数为三的 Calabi-Yau 几何，定义为
\begin{equation}
P_\Gamma(x,y)+P_{\Gamma'}(z,w)=0.\label{geomGG}
\end{equation}
例如，锥形奇点几何 \eqref{conifold} 对应于一对 $(A_1,A_1)$。稍微混用一下记号，将 IIB 型弦论放在由 $(G,G')$ 标记的几何上所得到的 $\cN{=}2$ 理论，也记为同一个符号 $(G,G')$。这类理论在文献 \cite{Cecotti:2010fi} 中引入并研究。显然有
\begin{equation}
(G,G')=(G',G).
\end{equation}
此外，只要写出几何 \eqref{geomGG}，下述等式便一目了然：
\begin{equation}
(D_4,A_3)=(E_6,A_2),\qquad(E_8,A_3)=(E_6,A_4).
\end{equation}

注意，也可以通过将 IIB 型理论放在 $\bC^2/\Gamma$ 上得到 $G_\Gamma$ 型六维 $\cN{=}(2,0)$ 理论。因此，第 \ref{sec:app6d} 小节中讨论的那类模型，也可以通过将 IIB 型理论放在复维数为三的 Calabi-Yau 几何上获得，该几何是形变后的 $\bC^2/\Gamma$ 奇点在紫外曲线 $C$ 上的纤维化。这里引入的理论 $(G,G')$ 则使用一对群，但当 $G'=A_k$ 时，它们可以与 $G$ 型六维构造联系起来：
\begin{equation}
(G,A_k)=Y_{2h^\vee+k+1}(G).
\end{equation}
结合式 \eqref{ady}，得到
\begin{equation}
AD_\text{pure}(G)=Y_{2h^\vee+2}(G)=(G,A_1).\label{ADG}
\end{equation}
进一步令 $G=A_{N-1}$，该等式变为
\begin{equation}
\begin{aligned}
AD_\text{pure}(A_{N-1})&=Y_{2N+2}(A_{N-1})=(A_{N-1},A_1)\\
&=(A_1,A_{N-1})=Y_{N+4}(A_1).
\end{aligned}
\end{equation}
再令 $G=A_1$，可见 $(A_1,A_1)$ 对应于纯 $\SU(2)$ 理论的最奇异点，即表示单个超多重态的单极子点。这与式 \eqref{conifold} 之后的陈述一致。

在式 \eqref{ADG} 中令 $G=D_N$，有
\begin{equation}
AD_\text{pure}(D_N)=Y_{4N-2}(D_N)=(D_N,A_1)
\end{equation}
但由式 \eqref{ADDD} 已知，这也等于 $X_{N+2}(A_1)$；正如式 \eqref{lda} 中提到的，它在文献 \cite{Cecotti:2013lda} 中记为 $D_N(A_1)$。

本节迄今引入的理论之间，还存在其他一些偶然等价。不过，或许等到对整个 $\cN{=}2$ 理论谱系有了更清晰的认识之后，再全面讨论这些重合会更合适。
